\begin{abstract}
\noindent\textbf{Question.} How much control over an in-plane elastic wave field can the
fiber orientation of a continuous-fiber composite exert, at fixed material and mass and
under the tow-curvature limit of automated placement? By what mechanism does it act, and
what bounds it?

\noindent\textbf{Hypothesis.} In a strongly anisotropic lamina, wave energy is skewed
toward the fiber direction: by up to $65^\circ$ for CFRP, so that $88\%$ of quasi-P wave
normals carry their energy within $20^\circ$ of the fibers. A fiber field should
therefore act as a designable field of energy streamlines. Its robust effect should be
\emph{redirection} of energy, which is broadband. It should vanish for isotropic material,
grow with anisotropy, and be limited by how sharply tows may turn.

\noindent\textbf{Analysis.} We solve the genuine two-dimensional orthotropic
time-harmonic problem in \texttt{dolfinx}, using an exact real four-field split and
automatic discrete adjoints, with a smooth radial-basis orientation field and a
pointwise tow-curvature constraint. We decompose every focus gain exactly into
redirection and field growth, sweep frequency on fixed designs, compute the energy-flux
field, re-optimize along an anisotropy family, and trace the curvature front.

\noindent\textbf{Results.}
\begin{itemize}[leftmargin=1.2em,itemsep=0pt,topsep=1pt]
\item An orientation-only lens's $\WGainForty\times$ single-frequency focus gain factorizes
  into $\WRedirForty\times$ redirection (target share $\WPhiZero\%\to\WPhiOne\%$) and a
  $\WGrowthForty\times$ resonant growth of the whole field. With a better absorber ($|R|\le\WReflThickMax$ instead of $\WReflMin$--$\WReflMax$) the re-optimized lens gives
  $\WAbsG\times=\WAbsP\times$ redirection $\times\,\WAbsGr\times$ growth: the split is
  robust, the redirection is absorber-dependent.
\item Only the redirection persists off-design: it stays $\ge\WRedirBandMin\times$ over
  $\omega\in[37.5,60]$ ($\ge\WAbsMinFortySixty\times$ over $[40,60]$ with the better
  absorber), while the gain collapses to $\WGainMinusFive\times$ at $-6\%$.
\item The energy flux follows the local fibers: its mean deviation from them is
  $\WFluxDev^\circ$ at the design frequency and at most $\WFluxDevThirtyFive^\circ$ across the
  band tested, and by only $\WCtrlDevMin$--$\WCtrlDevMax^\circ$ in curved fields that were
  never optimized.
\item Redirection is exactly absent for isotropic material and grows with anisotropy.
  It is unaffected while the minimum tow radius is at most about $0.7$ along-fiber
  wavelengths; beyond it the best design found declines gradually.
\item A cloak reduces the far-field mismatch $\WCloakPeak\times$ at its design frequency (a
  value tuned to its mesh: $\WCloakFineEval\times$ on a finer mesh, $\WCloakFineReopt\times$ when
  re-optimized there) but only $\WCloakMinus$--$\WCloakPlus\times$ at $\pm5\%$, as expected for a full cloak that
  relies on interference.
\end{itemize}

\noindent\textbf{Conclusion.} Fiber orientation steers in-plane elastic energy by channeling
it along the fibers. The transferable, broadband capability is redirection of a few
times. Large single-frequency gains of finite plates are mostly resonance and should not
be reported as steering. Designs should be optimized and judged over a band.
\end{abstract}

\noindent\textbf{Keywords:} elastic waves; fiber steering; anisotropic plates; energy channeling;
topology optimization; cloaking; automated fiber placement.

